\documentclass[10pt,a4paper]{amsart}
\usepackage{amsmath,amssymb,amsthm,graphicx}
\usepackage{verbatim}
\newenvironment{reference}{\comment}{\endcomment}
%\newenvironment{reference}{}{}
\newenvironment{slogan}{\comment}{\endcomment}
\newenvironment{history}{\comment}{\endcomment}
% For commutative diagrams we use Xy-pic
\usepackage[all]{xy}

% We use 2cell for 2-commutative diagrams.
\xyoption{2cell}
% Theorem environments.
%
\theoremstyle{plain}
\newtheorem{theorem}[subsection]{Theorem}
\newtheorem{proposition}[subsection]{Proposition}
\newtheorem{lemma}[subsection]{Lemma}

\theoremstyle{definition}
\newtheorem{definition}[subsection]{Definition}
\newtheorem{example}[subsection]{Example}
\newtheorem{exercise}[subsection]{Exercise}
\newtheorem{situation}[subsection]{Situation}

\theoremstyle{remark}
\newtheorem{remark}[subsection]{Remark}
\newtheorem{remarks}[subsection]{Remarks}

\numberwithin{equation}{subsection}

% Macros
%
\def\lim{\mathop{\mathrm{lim}}\nolimits}
\def\colim{\mathop{\mathrm{colim}}\nolimits}
\def\Spec{\mathop{\mathrm{Spec}}}
\def\Hom{\mathop{\mathrm{Hom}}\nolimits}
\def\Ext{\mathop{\mathrm{Ext}}\nolimits}
\def\SheafHom{\mathop{\mathcal{H}\!\mathit{om}}\nolimits}
\def\SheafExt{\mathop{\mathcal{E}\!\mathit{xt}}\nolimits}
\def\Sch{\mathit{Sch}}
\def\Mor{\mathop{\mathrm{Mor}}\nolimits}
\def\Ob{\mathop{\mathrm{Ob}}\nolimits}
\def\Sh{\mathop{\mathit{Sh}}\nolimits}
\def\NL{\mathop{N\!L}\nolimits}
\def\CH{\mathop{\mathrm{CH}}\nolimits}
\def\proetale{{pro\text{-}\acute{e}tale}}
\def\etale{{\acute{e}tale}}
\def\QCoh{\mathit{QCoh}}
\def\Ker{\mathop{\mathrm{Ker}}}
\def\Im{\mathop{\mathrm{Im}}}
\def\Coker{\mathop{\mathrm{Coker}}}
\def\Coim{\mathop{\mathrm{Coim}}}

% Boxtimes
%
\DeclareMathSymbol{\boxtimes}{\mathbin}{AMSa}{"02}

%
% Macros for moduli stacks/spaces
%
\def\QCohstack{\mathcal{QC}\!\mathit{oh}}
\def\Cohstack{\mathcal{C}\!\mathit{oh}}
\def\Spacesstack{\mathcal{S}\!\mathit{paces}}
\def\Quotfunctor{\mathrm{Quot}}
\def\Hilbfunctor{\mathrm{Hilb}}
\def\Curvesstack{\mathcal{C}\!\mathit{urves}}
\def\Polarizedstack{\mathcal{P}\!\mathit{olarized}}
\def\Complexesstack{\mathcal{C}\!\mathit{omplexes}}
% \Pic is the operator that assigns to X its picard group, usage \Pic(X)
% \Picardstack_{X/B} denotes the Picard stack of X over B
% \Picardfunctor_{X/B} denotes the Picard functor of X over B
\def\Pic{\mathop{\mathrm{Pic}}\nolimits}
\def\Picardstack{\mathcal{P}\!\mathit{ic}}
\def\Picardfunctor{\mathrm{Pic}}
\def\Deformationcategory{\mathcal{D}\!\mathit{ef}}
\usepackage{cleveref}
\crefname{originalreference}{original reference}{original references}
\Crefname{originalreference}{Original reference}{Original references}
\makeatletter
\newlabel{sheaves-lemma-stalk-pullback-modules}{{sheaves-lemma-stalk-pullback-modules}{source}{}{sourcepointer-1}{}}
\newlabel{sheaves-lemma-stalk-pullback-modules@cref}{{[originalreference][0][]sheaves-lemma-stalk-pullback-modules}{[1][0][]source}{}{}{}}
\makeatother
\crefname{originalreference}{original reference}{original references}
\Crefname{originalreference}{Original reference}{Original references}
\makeatletter
\newlabel{sheaves-lemma-stalks-closed-pushforward}{{sheaves-lemma-stalks-closed-pushforward}{source}{}{sourcepointer-2}{}}
\newlabel{sheaves-lemma-stalks-closed-pushforward@cref}{{[originalreference][0][]sheaves-lemma-stalks-closed-pushforward}{[1][0][]source}{}{}{}}
\makeatother
\title{Modules on a closed subspace are exactly the modules killed by the ideal of the immersion}
\author{The Stacks Project Authors}
\date{stacks-project:modules.tex:local-snapshot-20261012; GNU FDL 1.2 or later}
\begin{document}
\maketitle
\noindent\textit{Editorial note.} All delivered mathematical text is the author's own, copied from the pinned Stacks Project chapter source sources/stacks-modules/pinned-original-source.tex; no statement or proof step was rewritten, reordered or removed.\par\medskip
\noindent\textit{Editorial note.} Preamble provenance: the Stacks Project keeps one single shared preamble for every chapter; this card takes it directly from the staged sources/stacks-modules/preamble.tex as three verbatim spans (lines 8-12, 14-18 and 154-217), so no derived file was written and nothing was added or redefined. The wrapper class and the base packages are supplied by the assembler.\par\medskip
\noindent\textit{Editorial note.} Bibliography provenance: neither the counted statement nor its proof contains a citation (this lane rescans the whole span before writing the card), so the card carries no bibliography and no entry was reconstructed.\par\medskip
\noindent\textit{Editorial note.} Selection provenance: this card was opened under the file-based control contract (CONTROL.md), with the mechanical pre-card gate run on its statement and proof interval before the card was written.\par\medskip
\noindent\textit{Allowed standard prerequisites (editorial).}\begin{itemize}
\item The stalk of a pullback of modules and the description of the stalks of the direct image along a homeomorphism onto a closed subset. Both are invoked inside the counted proof (source lines 2192-2194 and 2208-2209, 2216-2217) as cross-chapter pointers. Allow the two canonical maps of the proof to be checked on stalks.
\item The kernel of the surjection of structure sheaves as an ideal sheaf, and the induced quotient description of the local rings of the subspace. The statement of the counted lemma and the proof at source lines 2217-2219. Identifies the local rings of the closed subspace with the quotients of those of the ambient space and makes the annihilation hypothesis usable.
\end{itemize}\medskip
\ifdefined\hypertarget\hypertarget{sourcepointer-1}{}\fi
\noindent\textit{Source reference.} Original reference sheaves-lemma-stalk-pullback-modules: sources/stacks-modules/pinned-original-source.tex, source lines 1418--1418. Navigation pointer only: the counted text refers to the Stacks item sheaves-lemma-stalk-pullback-modules, which belongs to another chapter of the project and is not part of this card. The pointer carries the target's own label as its number because this chapter's staged snapshot does not contain the other chapter's tag table; its content is not reproduced here.\par\medskip
\ifdefined\hypertarget\hypertarget{sourcepointer-2}{}\fi
\noindent\textit{Source reference.} Original reference sheaves-lemma-stalks-closed-pushforward: sources/stacks-modules/pinned-original-source.tex, source lines 787--787. Navigation pointer only: the counted text refers to the Stacks item sheaves-lemma-stalks-closed-pushforward, which belongs to another chapter of the project and is not part of this card. The pointer carries the target's own label as its number because this chapter's staged snapshot does not contain the other chapter's tag table; its content is not reproduced here.\par\medskip
\begin{lemma}
\label{lemma-i-star-equivalence}
Let $i : (Z, \mathcal{O}_Z) \to (X, \mathcal{O}_X)$ be a morphism
of ringed spaces. Assume $i$ is a homeomorphism onto a closed subset of $X$
and $i^\sharp : \mathcal{O}_X \to i_*\mathcal{O}_Z$ is surjective.
Denote $\mathcal{I} \subset \mathcal{O}_X$ the kernel of $i^\sharp$.
The functor
$$
i_* :
\textit{Mod}(\mathcal{O}_Z)
\longrightarrow
\textit{Mod}(\mathcal{O}_X)
$$
is exact, fully faithful, with essential image those
$\mathcal{O}_X$-modules $\mathcal{G}$ such that $\mathcal{I}\mathcal{G} = 0$.
\end{lemma}

\begin{proof}
We claim that for an $\mathcal{O}_Z$-module $\mathcal{F}$ the canonical map
$$
i^*i_*\mathcal{F} \longrightarrow \mathcal{F}
$$
is an isomorphism. We check this on stalks. Say $z \in Z$ and $x = i(z)$.
We have
$$
(i^*i_*\mathcal{F})_z =
(i_*\mathcal{F})_x \otimes_{\mathcal{O}_{X, x}} \mathcal{O}_{Z, z} =
\mathcal{F}_z \otimes_{\mathcal{O}_{X, x}} \mathcal{O}_{Z, z} =
\mathcal{F}_z
$$
by Sheaves, Lemma \ref{sheaves-lemma-stalk-pullback-modules},
the fact that $\mathcal{O}_{Z, z}$ is a quotient of $\mathcal{O}_{X, x}$, and
Sheaves, Lemma \ref{sheaves-lemma-stalks-closed-pushforward}.
It follows that $i_*$ is fully faithful.

\medskip\noindent
Let $\mathcal{G}$ be a $\mathcal{O}_X$-module with
$\mathcal{I}\mathcal{G} = 0$. We will prove the canonical map
$$
\mathcal{G} \longrightarrow i_*i^*\mathcal{G}
$$
is an isomorphism. This proves that $\mathcal{G} = i_*\mathcal{F}$
with $\mathcal{F} = i^*\mathcal{G}$ which finishes the proof.
We check the displayed map induces an isomorphism on stalks.
If $x \in X$, $x \not \in i(Z)$, then $\mathcal{G}_x = 0$
because $\mathcal{I}_x = \mathcal{O}_{X, x}$ in this
case. As above $(i_*i^*\mathcal{G})_x = 0$ by
Sheaves, Lemma \ref{sheaves-lemma-stalks-closed-pushforward}.
On the other hand, if $x \in Z$, then we obtain the map
$$
\mathcal{G}_x
\longrightarrow
\mathcal{G}_x \otimes_{\mathcal{O}_{X, x}} \mathcal{O}_{Z, x}
$$
by Sheaves, Lemmas \ref{sheaves-lemma-stalk-pullback-modules} and
\ref{sheaves-lemma-stalks-closed-pushforward}. This map is an isomorphism
because $\mathcal{O}_{Z, x} = \mathcal{O}_{X, x}/\mathcal{I}_x$
and because $\mathcal{G}_x$ is annihilated by $\mathcal{I}_x$ by assumption.
\end{proof}

\end{document}
